\section*{Capítulo \ref{ch:b2:diffcalc} --- Cálculo diferencial}

\begin{solution}{exo:b2:diffcalc:1}
$J_f = \begin{pmatrix} 2x & -2y\\ 2y & 2x\end{pmatrix}$, $\det J_f
= 4(x^2 + y^2)$: invertível fora da origem. (Essa $f$ é $z
\mapsto z^2$ disfarçada de complexa.)

$J_\Phi = \begin{pmatrix} \cos\theta & -r\sin\theta & 0\\
\sin\theta & r\cos\theta & 0\\ 0 & 0 & 1\end{pmatrix}$, $\det = r$:
invertível para $r \neq 0$ (coordenadas cilíndricas).
\end{solution}

\begin{solution}{exo:b2:diffcalc:2}
Para $v = (a, b) \neq 0$: $\frac{f(tv) - 0}{t} = \frac{t^3a^3}{t\cdot
t^2(a^2+b^2)} = \frac{a^3}{a^2 + b^2}$: toda derivada
direcional existe, com valor $D_v = \frac{a^3}{a^2+b^2}$. Mas $v
\mapsto D_v$ não é linear ($D_{(1,0)} = 1$, $D_{(0,1)} = 0$,
$D_{(1,1)} = \frac12 \neq 1$): nenhuma aplicação linear pode produzir esses
valores, logo $f$ não é \omterm{def:b2:diffcalc:differential}{diferenciável} em $0$ (a \omterm{def:b2:diffcalc:differential}{diferencial} teria
de ser $v \mapsto D_v$).
\end{solution}

\begin{solution}{exo:b2:diffcalc:3}
$f = x^3 + y^3 - 3xy$: pontos críticos $(0,0)$ e $(1,1)$ (cálculo do primeiro
ano). Hessianas: $H = \begin{pmatrix} 6x & -3\\ -3 &
6y\end{pmatrix}$. Em $(0,0)$: sinais próprios mistos ($\det = -9 < 0$):
sela. Em $(1,1)$: $\det = 27 > 0$, traço $> 0$: positiva
definida, mínimo local estrito --- agora justificado pelo
\cref{thm:b2:diffcalc:extrema} em vez de decretado.

$g = x^4 + y^4 - 2(x-y)^2$: $\nabla g = (4x^3 - 4(x - y),\; 4y^3 +
4(x-y))$; pontos críticos $(0,0)$, $(\sqrt2, -\sqrt2)$,
$(-\sqrt2, \sqrt2)$ (primeiro ano). Em $(\pm\sqrt2, \mp\sqrt2)$: $H =
\begin{pmatrix} 12\cdot2 - 4 & 4\\ 4 & 20\end{pmatrix} =
\begin{pmatrix} 20 & 4\\ 4 & 20\end{pmatrix}$: positiva definida
(diagonalmente dominante; \omterm{def:b2:reduction:eigen}{autovalores} $24, 16$): mínimos locais estritos.
Em $(0,0)$: $H = \begin{pmatrix} -4 & 4\\ 4 & -4\end{pmatrix}$,
singular negativa semidefinida: o teste se cala; o estudo
direcional ($g(x,x) = 2x^4 > 0$, $g(x,-x) = 2x^4 - 8x^2 < 0$ pequeno)
mostra um ponto do tipo sela --- sem extremo.
\end{solution}

\begin{solution}{exo:b2:diffcalc:4}
Derive $t \mapsto f(tx)$ em $t = 1$: pela regra da cadeia,
$\langle \nabla f(x), x\rangle$; por homogeneidade, a mesma função
é $t^pf(x)$, de derivada $pf(x)$ em $t = 1$: a identidade de Euler.

Recíproca: fixe $x \neq 0$ e ponha $\varphi(t) = f(tx) - t^p f(x)$
em $t > 0$. Então $\varphi'(t) = \langle\nabla f(tx), x\rangle -
pt^{p-1}f(x) = \frac1t\bigl(\langle \nabla f(tx), tx\rangle -
p\,t^pf(x)\bigr)$. A hipótese --- a identidade de Euler no ponto
$tx$ --- avalia o colchete como $p\,f(tx) - p\,t^pf(x) =
p\,\varphi(t)$. Logo $\varphi' = \frac{p}{t}\varphi$ com $\varphi(1)
= 0$: a única solução da EDO linear é $\varphi \equiv 0$ (unicidade do
primeiro ano), isto é, $f(tx) = t^pf(x)$.
\end{solution}

\begin{solution}{exo:b2:diffcalc:5}
Desenvolvendo $f(x + h) - f(x) = \langle Ax - b, h\rangle +
\frac12\langle Ah, h\rangle$ (simetria de $A$): $\nabla f(x) = Ax -
b$ e $H_f = A$ em toda parte. $f$ é convexa se e somente se $A$ é positiva
semidefinida (o teste da hessiana, global aqui pois $H$ é constante:
a fórmula de Taylor de segunda ordem é exata). Se $A$ é positiva
definida: o único ponto crítico é $x^* = A^{-1}b$, e $f(x^*
+ h) - f(x^*) = \frac12\langle Ah, h\rangle \geq
\frac{\lambda_{\min}}{2}\norm h^2 > 0$ para $h \neq 0$: mínimo
global estrito.
\end{solution}

\begin{solution}{exo:b2:diffcalc:6}
Como $\nabla g(a) \neq 0$, alguma parcial, digamos $\frac{\partial
g}{\partial y}(a) \neq 0$: o teorema da função implícita
(companheiro do~\cref{thm:b2:diffcalc:inverse}) parametriza $\{g =
0\}$ perto de $a = (a_1, a_2)$ como $\gamma(t) = (t, y(t))$ com $y$
$C^1$, $y'(t) = -\frac{\partial_x g}{\partial_y g}(\gamma(t))$
(derive $g(t, y(t)) = 0$). A função de uma variável $t
\mapsto f(\gamma(t))$ tem extremo local em $t = a_1$:
\[
0 = \frac{\dd}{\dd t}f(\gamma(t))\Big|_{a_1}
= \partial_x f(a) + \partial_y f(a)\,y'(a_1)
= \partial_x f(a) - \partial_yf(a)\frac{\partial_x
g(a)}{\partial_y g(a)} :
\]
os vetores $\nabla f(a)$ e $\nabla g(a)$ têm coordenadas
proporcionais: $\nabla f(a) = \lambda \nabla g(a)$ com $\lambda =
\frac{\partial_y f(a)}{\partial_y g(a)}$.

Aplicação: no círculo, $\nabla(xy) = (y, x)$ paralelo a
$(2x, 2y)$ força $y^2 = x^2$; com a restrição, os candidatos
são $\pm\bigl(\tfrac{1}{\sqrt2}, \tfrac{1}{\sqrt2}\bigr)$ (valor
$\frac12$) e $\pm\bigl(\tfrac{1}{\sqrt2},
-\tfrac{1}{\sqrt2}\bigr)$ (valor $-\frac12$): máximo $\frac12$, mínimo
$-\frac12$ (atingidos: o círculo é \omterm{def:b2:metric:compact}{compacto}).
\end{solution}

\begin{solution}{exo:b2:diffcalc:7}
$\nabla f = (2x - y + 1,\; 2y - x - 1) = 0$: resolvendo, $x =
-\frac13$, $y = \frac13$. Hessiana $\begin{pmatrix} 2 & -1\\ -1 &
2\end{pmatrix}$, positiva definida: mínimo global da quadrática
$f$, valor $f\bigl(-\frac13, \frac13\bigr) = -\frac13$.

No triângulo $T$: o ponto crítico interior $(-\frac13,
\frac13) \notin T$ (negativo, $x$). Lados: em $y = 0$, $x \in
\intcc{0}{1}$: $f = x^2 + x$, crescente: extremos $0$ e $2$. Em
$x = 0$: $f = y^2 - y$, mínimo $-\frac14$ em $y = \frac12$, valores
$0$ e $0$ nas pontas. Em $x + y = 1$: substitua $y = 1 - x$,
$f = x^2 + (1-x)^2 - x(1-x) + x - (1-x) = 3x^2 - x$; em
$\intcc{0}{1}$: mínimo $-\frac{1}{12}$ em $x = \frac16$, valores
$0$ (em $x=0$) e $2$ (em $x=1$). Vértices: $f(0,0) = 0$, $f(1,0)
= 2$, $f(0,1) = 0$. Global em $T$: mínimo $-\frac14$ em $(0,
\frac12)$, máximo $2$ em $(1, 0)$.
\end{solution}

\begin{solution}{exo:b2:diffcalc:8}
$J_f = \begin{pmatrix} 1 & 2y\\ 2x & 1\end{pmatrix}$, $\det J_f =
1 - 4xy$. Em $0$: $\det = 1 \neq 0$: difeomorfismo local
(\cref{thm:b2:diffcalc:inverse}), com
\[
\dd(f^{-1})_{(0,0)} = (J_f(0))^{-1} = I_2 .
\]
Invertibilidade numa bola: é preciso $4\abs{xy} < 1$ em toda parte; em
$\norm{(x,y)}_2 < r$, $\abs{xy} \leq \frac{x^2 + y^2}{2} <
\frac{r^2}{2}$, logo $r = \frac{1}{\sqrt2}$ serve; e é o maior:
em $(x, y) = \bigl(\tfrac12, \tfrac12\bigr)$, de \omterm{def:b2:nvs:norm}{norma}
$\frac{1}{\sqrt2}$, $\det J_f = 0$.
\end{solution}

\begin{solution}{exo:b2:diffcalc:9}
$f(t) = (\cos t, \sin t)$: $f(0) = f(2\pi) = (1, 0)$, e no entanto $f'(t) =
(-\sin t, \cos t)$ tem \omterm{def:b2:nvs:norm}{norma} $1$, nunca nula: não há ponto em que a
derivada se anule --- Rolle não tem análogo vetorial. A desigualdade
do valor médio vale confortavelmente: $\norm{f(2\pi) - f(0)} = 0 \leq
2\pi \cdot \sup\norm{f'} = 2\pi$.
\end{solution}

\begin{solution}{exo:b2:diffcalc:10}
$f(x+h) - f(x) = 2\langle x, h\rangle + \norm h^2$: $\dd f_x =
2\langle x, \cdot\rangle$, $\nabla f(x) = 2x$, hessiana $2I$
(constante). Para $g$:
\[
g(x + h) - g(x) = \langle Ax, h\rangle + \langle Ah, x\rangle
+ \langle Ah, h\rangle
= \bigl\langle (A + A^{\mathsf T})x,\ h\bigr\rangle +
O(\norm h^2),
\]
logo $\nabla g(x) = (A + A^{\mathsf T})x$ e $H_g = A +
A^{\mathsf T}$, constante. Pelo~\cref{exo:b2:diffcalc:11}, $g$ é
convexa se e somente se $A + A^{\mathsf T}$ é positiva semidefinida ---
só a parte \omterm{def:b2:quadratic:adjoint}{simétrica} de $A$ importa, como de fato $g(x) =
\langle \frac{A + A^{\mathsf T}}2 x, x\rangle$.
\end{solution}

\begin{solution}{exo:b2:diffcalc:11}
$f$ é convexa se e somente se sua restrição a todo segmento é convexa,
isto é, se e somente se toda $\varphi(t) = f(a + tv)$ é convexa. Pela regra
da cadeia, $\varphi''(t) = \langle H_{a+tv}\,v,\ v\rangle$.

Se todas as hessianas são positivas semidefinidas: $\varphi'' \geq 0$,
de modo que cada $\varphi$ é convexa (volume do primeiro ano de graduação) e $f$ é convexa.
Reciprocamente, se $f$ é convexa, cada $\varphi$ é convexa, logo
$\varphi''(0) \geq 0$: $\langle H_a v, v\rangle \geq 0$ para
todo $a$ e toda direção $v$: todas as hessianas são positivas
semidefinidas.
\end{solution}

\begin{solution}{exo:b2:diffcalc:12}
\begin{enumerate}
  \item Pela desigualdade do valor médio
        (\cref{thm:b2:diffcalc:mvi}) aplicada a $g$:
        $\norm{g(x) - g(y)} \leq k\norm{x-y}$, logo
        \[
        \norm{f(x) - f(y)} \geq \norm{x - y} - \norm{g(x) -
        g(y)} \geq (1 - k)\norm{x - y} :
        \]
        $f$ é injetiva e $f^{-1}$ (definida na imagem) é
        $\frac{1}{1-k}$-lipschitziana.
  \item Fixe $y$; $T(x) = y - g(x)$ satisfaz $\norm{T(x) -
        T(x')} = \norm{g(x') - g(x)} \leq k\norm{x - x'}$: uma
        contração do espaço \omterm{def:b2:metric:complete}{completo} $\R^n$. Banach
        (\cref{thm:b2:metric:banach}) dá um ponto fixo $x^*
        = y - g(x^*)$, isto é, $f(x^*) = y$: $f$ é sobrejetiva.
  \item $f$ é assim uma bijeção de $\R^n$ com
        inversa $\frac{1}{1-k}$-lipschitziana: um teorema de inversão
        global, em que a pequenez de $\dd g$ em toda parte
        substitui a hipótese de invertibilidade local do
        \cref{thm:b2:diffcalc:inverse}.
\end{enumerate}
\end{solution}

\begin{solution}{pb:b2:diffcalc:1}
\textbf{1.} $\norm{ABx} \leq \vertiii A\norm{Bx} \leq
\vertiii A\vertiii B\norm x$: tome o sup sobre $\norm x = 1$.
Os produtos de matrizes, $\det$ e $\operatorname{tr}$ são funções
polinomiais das entradas, logo \omterm{def:b2:metric:continuity}{contínuas} ($\mathcal M_n(\R)
\simeq \R^{n^2}$, todas as \omterm{def:b2:nvs:equivalent}{normas equivalentes}:
\cref{thm:b2:nvs:finitedim}).

\textbf{2.} $\sum\vertiii{X^k} \leq \sum\vertiii X^k <
\infty$: a série converge \omterm{def:b2:series:def}{absolutamente}, logo converge
(\cref{thm:b2:nvs:absoluteconvergence}). De $(I -
X)\sum_{k\leq N}X^k = I - X^{N+1} \to I$: a soma é $(I -
X)^{-1}$. \omterm{def:b2:nvs:norm}{Norma}: $\leq \sum\vertiii X^k = \frac{1}{1 -
\vertiii X}$; e
\[
(I - X)^{-1} - I - X = \sum_{k\geq2}X^k = X^2(I - X)^{-1},
\qquad
\vertiii{X^2(I-X)^{-1}} \leq
\frac{\vertiii X^2}{1 - \vertiii X} = O(\vertiii X^2).
\]

\textbf{3.} $A + H = A(I + A^{-1}H)$ com $\vertiii{A^{-1}H}
\leq \vertiii{A^{-1}}\vertiii H < 1$: invertível pela questão 2:
a bola \omterm{def:b2:metric:topology}{aberta} de raio $\vertiii{A^{-1}}^{-1}$ em torno de $A$ está
em $GL_n(\R)$. Densidade: $\det(A + \varepsilon I)$ é um polinômio de grau $n$
em $\varepsilon$ com coeficiente dominante $1$: ele tem
finitas raízes, de modo que existem $\varepsilon_k \to 0$
com $A + \varepsilon_kI$ invertível, convergindo para $A$.

\textbf{4.} Para $\vertiii H < \frac{1}{2\vertiii{A^{-1}}}$:
\[
(A + H)^{-1} - A^{-1}
= \bigl[(I + A^{-1}H)^{-1} - I\bigr]A^{-1},
\]
de \omterm{def:b2:nvs:norm}{norma} no máximo $\frac{\vertiii{A^{-1}H}}{1 -
\vertiii{A^{-1}H}}\,\vertiii{A^{-1}} \leq
2\vertiii{A^{-1}}^2\vertiii H \to 0$: $\Phi$ é \omterm{def:b2:metric:continuity}{contínua} em
todo $A \in GL_n(\R)$.

\textbf{5.} $(A+H)^2 = A^2 + AH + HA + H^2$: a aplicação $H
\mapsto AH + HA$ é linear e o erro $H^2$ é
$O(\vertiii H^2)$. Desenvolvendo $(A + H)^k$ e ordenando pelo
número de fatores $H$: a parte linear é
$\sum_{i=0}^{k-1}A^iHA^{k-1-i}$, e os termos com $\geq 2$
fatores $H$ são majorados por $\binom k2$ produtos de \omterm{def:b2:nvs:norm}{norma}
de escala $\leq \vertiii A^{k-2}\vertiii H^2$: $O(\vertiii H^2)$.
Não se pode colapsar a soma em $kA^{k-1}H$ porque $H$ e $A$
não precisam comutar --- a soma é a derivada não comutativa
correta.

\textbf{6.} Para $H$ pequena:
\[
(A+H)^{-1} = (I + A^{-1}H)^{-1}A^{-1}
= \bigl(I - A^{-1}H + O(\vertiii H^2)\bigr)A^{-1}
= A^{-1} - A^{-1}HA^{-1} + O(\vertiii H^2) :
\]
$\dd\Phi_A(H) = -A^{-1}HA^{-1}$, linear em $H$. Para $n = 1$:
$\dd(1/a)(h) = -h/a^2$, a derivada familiar.

\textbf{7.} Regra da cadeia ao longo da curva:
$\bigl(A(t)^{-1}\bigr)' = \dd\Phi_{A(t)}(A'(t)) =
-A(t)^{-1}A'(t)A(t)^{-1}$. Em $A(t) = I + tB$, $t = 0$: $(I +
tB)^{-1} = I - tB + O(t^2)$.

\textbf{8.} Pela questão 5 e pela invariância \omterm{def:b2:structures:generated}{cíclica} do traço:
\[
\dd f_A(H) = \operatorname{tr}\Bigl(\sum_{i=0}^{k-1}
A^iHA^{k-1-i}\Bigr) = k\operatorname{tr}\bigl(A^{k-1}H\bigr) .
\]
Contra o produto de Frobenius, $\dd f_A(H) =
\operatorname{tr}\bigl((\nabla f)^{\mathsf T}H\bigr)$ exige
$(\nabla f)^{\mathsf T} = kA^{k-1}$: $\nabla f(A) =
k\,(A^{k-1})^{\mathsf T}$.

\textbf{9.} $f(A + H) - f(A) = 2\operatorname{tr}
(A^{\mathsf T}H) + \operatorname{tr}(H^{\mathsf T}H)$: a
\omterm{def:b2:diffcalc:differential}{diferencial} é $H \mapsto 2\operatorname{tr}(A^{\mathsf T}H) =
2\langle A, H\rangle$, logo $\nabla f(A) = 2A$; o termo de
segunda ordem é exatamente $\norm H_F^2$: a hessiana é o dobro da
\omterm{def:b2:quadratic:def}{forma quadrática} identidade, positiva definida e constante, de modo que
$\norm\cdot_F^2$ é estritamente convexa (a fórmula de Taylor é
exata aqui).

\textbf{10.} Pela multilinearidade nas colunas, $\det(I + H) =
\sum_{S\subseteq\{1,\dots,n\}}\det(M_S)$, em que $M_S$ tem coluna
$h_j$ para $j \in S$ e $e_j$ caso contrário. $S = \varnothing$
dá $1$; $S = \{j\}$ dá o \omterm{def:b2:linalg:det}{determinante} de $I$ com a coluna
$j$ substituída por $h_j$, a saber sua $j$-ésima entrada $h_{jj}$,
somando $\operatorname{tr}H$; cada termo com $\abs S \geq 2$
é um \omterm{def:b2:linalg:det}{determinante} com ao menos duas colunas de tamanho
$O(\vertiii H)$, logo $O(\vertiii H^2)$ (aplicações multilineares num espaço de
dimensão finita são limitadas), e há um número finito
delas. Assim $\det(I + H) = 1 + \operatorname{tr}H +
O(\vertiii H^2)$: $\dd(\det)_I = \operatorname{tr}$.

\textbf{11.} $\det(A + H) = \det A\,\det(I + A^{-1}H) =
\det A\,\bigl(1 + \operatorname{tr}(A^{-1}H) +
O(\vertiii H^2)\bigr)$: a \omterm{def:b2:diffcalc:differential}{diferencial} é $H \mapsto
\det(A)\operatorname{tr}(A^{-1}H)$.

\textbf{12.} Desenvolvimento de Laplace ao longo da linha $i$: $\det A =
\sum_j a_{ij}C_{ij}$, e os cofatores $C_{ij}$ não envolvem
a linha $i$: $\frac{\partial\det}{\partial a_{ij}} = C_{ij}$. Logo
\[
\dd(\det)_A(H) = \sum_{i,j}C_{ij}h_{ij}
= \operatorname{tr}\bigl(\operatorname{com}(A)^{\mathsf T}
H\bigr)
= \operatorname{tr}\bigl(\operatorname{adj}(A)H\bigr),
\qquad
\nabla(\det)(A) = \operatorname{com}(A) .
\]
Para $A$ invertível, $\operatorname{adj}A = \det(A)A^{-1}$
recupera a questão 11. Ao longo de uma curva $C^1$, a regra da cadeia lê-se
$(\det A(t))' = \operatorname{tr}(\operatorname{adj}
(A(t))\,A'(t))$: a fórmula de Jacobi.

\textbf{13.} Com $A' = MA$ e
$\operatorname{adj}(A)A = \det(A)I$:
\[
(\det A)' = \operatorname{tr}\bigl(\operatorname{adj}(A)MA
\bigr)
= \operatorname{tr}\bigl(A\operatorname{adj}(A)M\bigr)
= \det A\;\operatorname{tr}M
\]
(ciclicidade; $A\operatorname{adj}A = \det(A) I$ também). A
EDO linear escalar $y' = \operatorname{tr}(M(t))\,y$ tem
solução única $y(t) = y(0)\exp\bigl(\int_0^t
\operatorname{tr}M\bigr)$ (primeiro ano): a fórmula de Liouville.

\textbf{14.} Em $A \in SL_n(\R)$, tome $H = \frac1nA$:
$\dd(\det)_A\bigl(\tfrac1nA\bigr) = \frac1n\det(A)
\operatorname{tr}(A^{-1}A) = \frac1n\cdot1\cdot n = 1 \neq 0$:
a \omterm{def:b2:diffcalc:differential}{diferencial} é uma forma linear não nula, logo sobrejetiva
sobre $\R$ em todo ponto do conjunto de nível: $SL_n(\R)$ é um
conjunto de nível suave (de dimensão $n^2 - 1$).

\textbf{15.} $\sum_k\vertiii{A^k/k!} \leq
\sum\vertiii A^k/k! = \eu^{\vertiii A}$: convergência absoluta
(o argumento de \omterm{def:b2:metric:complete}{completude} da questão 2), com \omterm{def:b2:funcseq:series}{convergência normal}
em toda bola $\vertiii A \leq R$ (majore $R^k/k!$ independentemente
de $A$). Cada soma parcial é \omterm{def:b2:metric:continuity}{contínua} (polinomial); o
limite uniforme em bolas é \omterm{def:b2:metric:continuity}{contínuo}: $A \mapsto \eu^A$ é
\omterm{def:b2:metric:continuity}{contínua}, com $\vertiii{\eu^A} \leq \eu^{\vertiii A}$.

\textbf{16.} As duas séries convergem \omterm{def:b2:series:def}{absolutamente}, de modo que o \omterm{thm:b2:series:fubini}{produto
de Cauchy} é legítimo (\cref{thm:b2:series:fubini}):
\[
\eu^A\eu^B = \sum_{n\geq0}\frac{1}{n!}\sum_{k=0}^n\binom
nkA^kB^{n-k}
= \sum_{n\geq0}\frac{(A+B)^n}{n!} = \eu^{A+B},
\]
exigindo a identidade binomial que $AB = BA$. Com $B = -A$:
$\eu^A\eu^{-A} = \eu^0 = I$: toda $\eu^A \in GL_n(\R)$.

\textbf{17.} A série $\sum t^kA^k/k!$ e sua série derivada
$\sum t^{k-1}A^k/(k-1)! = A\sum t^{k-1}A^{k-1}/(k-1)!$ convergem
\omterm{def:b2:funcseq:series}{normalmente} em todo segmento $\abs t \leq T$ (cotas
$T^k\vertiii A^k/k!$): o teorema de derivação para séries
(\cref{thm:b2:funcseq:seriestransfer}, aplicado entrada a entrada)
dá $\frac{\dd}{\dd t}\eu^{tA} = A\eu^{tA}$; fatorar $A$
à direita, em vez disso, dá $\eu^{tA}A$. Iterando: $C^\infty$.

\textbf{18.} $A(t) = \eu^{tA}$ satisfaz $A'(t) = A\,A(t)$:
a fórmula de Liouville (questão 13) com $M = A$ constante dá
$\det\eu^{tA} = \eu^{t\operatorname{tr}A}$ (valor $1$ em $t =
0$); em $t = 1$, $\det\eu^A = \eu^{\operatorname{tr}A}$.
Verificações: $n = 1$ é a própria exponencial; uma $A$ nilpotente tem
$\operatorname{tr}A = 0$ e $\eu^A$ unipotente de \omterm{def:b2:linalg:det}{determinante}
$1$; $\operatorname{tr}A = 0$ dá $\det\eu^A = 1$: as
matrizes de traço nulo são enviadas em $SL_n(\R)$.

\textbf{19.} $\eu^H - I - H = \sum_{k\geq2}H^k/k!$, de \omterm{def:b2:nvs:norm}{norma}
$\leq \vertiii H^2\eu^{\vertiii H} = O(\vertiii H^2)$:
$\dd(\exp)_0 = \mathrm{id}$, invertível. Além disso, $\exp$ é
$C^1$: pela questão 5, a candidata a \omterm{def:b2:diffcalc:differential}{diferencial} $H \mapsto
\sum_k\frac1{k!}\sum_iA^iHA^{k-1-i}$ é uma série \omterm{def:b2:funcseq:series}{normalmente} convergente
de aplicações lineares dependendo \omterm{def:b2:metric:continuity}{continuamente} de $A$ (cotas
$\vertiii A^{k-1}/(k-1)!$ em bolas), de modo que as parciais existem e
são \omterm{def:b2:metric:continuity}{contínuas} (\cref{thm:b2:diffcalc:c1} e o teorema de
transferência para séries). O teorema da função inversa
(\cref{thm:b2:diffcalc:inverse}) aplica-se em $0$: $\exp$ é um
difeomorfismo $C^1$ de uma vizinhança de $0$ sobre uma
vizinhança de $I$ --- matrizes perto de $I$ têm logaritmos.

\textbf{20.} Para $S = PDP^{\mathsf T}$ \omterm{def:b2:quadratic:adjoint}{simétrica} (teorema
espectral): $\eu^S = P\eu^DP^{\mathsf T}$ é \omterm{def:b2:quadratic:adjoint}{simétrica} com
\omterm{def:b2:reduction:eigen}{autovalores} $\eu^{\lambda_i} > 0$: positiva definida.
Sobrejetividade: uma $Q =
P\operatorname{diag}(\mu_i)P^{\mathsf T}$ positiva definida ($\mu_i > 0$) é
$\eu^S$ para $S = P\operatorname{diag}(\ln\mu_i)P^{\mathsf T}$.
Injetividade: $\eu^S$ determina seus \omterm{def:b2:reduction:eigen}{autoespaços}, que são
exatamente os de $S$ (em cada \omterm{def:b2:reduction:eigen}{autoespaço} de $S$ associado a $\lambda$,
$\eu^S$ age como $\eu^\lambda$; $\lambda$ distintos dão
$\eu^\lambda$ distintos), e tomar $\ln$ dos \omterm{def:b2:reduction:eigen}{autovalores}
recupera $S$. Assim $\exp$ é uma bijeção das matrizes \omterm{def:b2:quadratic:adjoint}{simétricas}
sobre as positivas definidas.

\textbf{21.} $F(A + H) = A^{\mathsf T}A + A^{\mathsf T}H +
H^{\mathsf T}A + H^{\mathsf T}H$: $\dd F_A(H) =
A^{\mathsf T}H + H^{\mathsf T}A$ (valores em $S_n$; erro
$O(\vertiii H^2)$). Em $A \in O_n$ e para $S \in S_n$, a
escolha $H = \frac12AS$ dá
\[
A^{\mathsf T}\cdot\tfrac12AS + \tfrac12(AS)^{\mathsf T}A
= \tfrac12 S + \tfrac12 S^{\mathsf T} = S :
\]
sobrejetiva. $O_n = F^{-1}(I)$ é um conjunto de nível suave de
dimensão $n^2 - \dim S_n = \frac{n(n-1)}2$.

\textbf{22.} Derivando $A(t)^{\mathsf T}A(t) = I$ em $t
= 0$ (com $A(0) = I$): $A'(0)^{\mathsf T} + A'(0) = 0$:
antissimétrica. Reciprocamente, para $K^{\mathsf T} = -K$:
$(\eu^{tK})^{\mathsf T}\eu^{tK} = \eu^{tK^{\mathsf T}}
\eu^{tK} = \eu^{-tK}\eu^{tK} = I$ (transponha a série
termo a termo; os expoentes comutam): a curva permanece em $O_n$,
com velocidade $K$ em $t = 0$. O espaço tangente em $I$ $=$ o das
matrizes antissimétricas, da dimensão esperada
$\frac{n(n-1)}2$.

\textbf{23.} $\operatorname{tr}K = 0$ para $K$ antissimétrica,
logo $\det\eu^K = \eu^0 = 1$ (questão 18): a exponencial
cai em $SO_n$. Para $n = 2$: $J^2 = -I$, logo $J^{2m} =
(-1)^mI$, $J^{2m+1} = (-1)^mJ$, e
\[
\eu^{\theta J}
= \Bigl(\sum_m\frac{(-1)^m\theta^{2m}}{(2m)!}\Bigr)I
+ \Bigl(\sum_m\frac{(-1)^m\theta^{2m+1}}{(2m+1)!}\Bigr)J
= \cos\theta\,I + \sin\theta\,J ,
\]
a rotação de $\theta$: as séries do seno e do cosseno vivem
dentro da \omterm{ex:b2:nvs:matrixexp}{exponencial de matriz}.

\textbf{24.} Para $A, B$ invertível: $\operatorname{adj}(AB) =
\det(AB)(AB)^{-1} = \det(B)\det(A)B^{-1}A^{-1} =
\operatorname{adj}(B)\operatorname{adj}(A)$. Os dois lados da
identidade são aplicações polinomiais (logo \omterm{def:b2:metric:continuity}{contínuas}) das entradas
de $(A, B)$; eles coincidem no subconjunto denso $GL_n\times GL_n$
de $\mathcal M_n\times\mathcal M_n$ (questão 3: aproxime
cada fator), logo coincidem em toda parte.

\textbf{25.} (i) A Parte I rodou sobre a \omterm{def:b2:metric:complete}{completude} dos espaços normados
de dimensão finita (as séries \omterm{def:b2:series:def}{absolutamente} convergentes convergem), a Parte IV
sobre a mesma coisa mais o teorema espectral na questão 20, e o logaritmo local
da Parte V sobre o teorema da função inversa. (ii) O
capítulo de equações diferenciais apoia-se na fórmula de Liouville
(questão 13) para o wronskiano de sistemas lineares, e em
$\frac{\dd}{\dd t}\eu^{tA} = A\eu^{tA}$ (questão 17), que é
o enunciado de que $\eu^{tA}$ resolve $X' = AX$. (iii)
$\dd(\det)_I = \operatorname{tr}$ diz que o traço é a
taxa infinitesimal de variação de volume, e $\det\eu^A =
\eu^{\operatorname{tr}A}$ integra globalmente esse enunciado.
(iv) O volume do terceiro ano de graduação nomeia as estruturas: $O_n$ e
$SL_n(\R)$ são grupos de Lie, seus espaços tangentes em $I$
(matrizes antissimétricas e de traço nulo) são \omterm{def:b2:structures:algebra}{álgebras} de Lie, e
$\exp$ é a ponte entre eles.
\end{solution}
